\capitulo{7}{Conclusiones y Líneas de trabajo futuras}

\section{Conclusiones}

El desarrollo de este Trabajo Fin de Grado ha permitido diseñar e implementar un sistema SIEM ligero orientado a la monitorización de eventos de seguridad dentro de un entorno simulado de alta seguridad basado en un polvorín militar.

A lo largo del proyecto se han aplicado conocimientos relacionados con programación, bases de datos, desarrollo de APIs, monitorización de eventos y diseño de interfaces web. Asimismo, se ha podido comprobar la importancia que tienen los sistemas SIEM dentro de los entornos donde resulta necesario supervisar de forma continua diferentes elementos de seguridad.

La arquitectura modular adoptada ha permitido desarrollar una solución flexible y fácilmente ampliable, facilitando la incorporación de nuevas funcionalidades y fuentes de información en futuras versiones.

\section{Objetivos alcanzados}

Los objetivos planteados al inicio del proyecto han sido alcanzados de forma satisfactoria.

Se ha desarrollado un sistema capaz de procesar eventos procedentes de diferentes dispositivos simulados, almacenarlos en una base de datos y aplicar mecanismos básicos de correlación y detección de anomalías.

Además, se ha implementado una API REST que permite consultar la información generada por el sistema y un dashboard web que facilita la visualización de eventos, alertas, anomalías y estado de los sensores.

El resultado obtenido constituye una solución funcional que reproduce los principios básicos de funcionamiento presentes en sistemas SIEM utilizados en entornos reales.

\section{Líneas de trabajo futuras}

Aunque el sistema desarrollado cumple los objetivos establecidos inicialmente, existen diversas líneas de mejora que podrían incorporarse en futuras versiones.

Entre las posibles ampliaciones destacan la integración con dispositivos reales de monitorización, la incorporación de algoritmos avanzados de aprendizaje automático para la detección de anomalías, la implementación de sistemas reales de notificación mediante correo electrónico o SMS y la integración con plataformas de monitorización externas.

También podría ampliarse la capacidad de correlación mediante reglas más complejas y mecanismos de análisis de comportamiento capaces de detectar patrones de actividad sospechosos de forma automática.

Por último, sería posible sustituir la base de datos SQLite por soluciones más escalables orientadas a entornos con un volumen elevado de eventos, permitiendo así utilizar el sistema en escenarios de mayor tamaño y complejidad.

En definitiva, el proyecto ha permitido desarrollar una solución funcional que integra los principales componentes presentes en un sistema SIEM moderno, proporcionando una visión práctica de los procesos de monitorización, correlación y análisis de eventos de seguridad.
